%
% File acl2015.tex
%
% Contact: car@ir.hit.edu.cn, gdzhou@suda.edu.cn
%%
%% Based on the style files for ACL-2014, which were, in turn,
%% Based on the style files for ACL-2013, which were, in turn,
%% Based on the style files for ACL-2012, which were, in turn,
%% based on the style files for ACL-2011, which were, in turn, 
%% based on the style files for ACL-2010, which were, in turn, 
%% based on the style files for ACL-IJCNLP-2009, which were, in turn,
%% based on the style files for EACL-2009 and IJCNLP-2008...

%% Based on the style files for EACL 2006 by 
%%e.agirre@ehu.es or Sergi.Balari@uab.es
%% and that of ACL 08 by Joakim Nivre and Noah Smith

\documentclass[11pt]{article}
\usepackage{acl2015}
\usepackage{times}
\usepackage{url}
\usepackage{latexsym}
\usepackage{natbib}
\usepackage{color}
\definecolor{citecolor}{RGB}{34,139,34}
\usepackage[pagebackref=true,breaklinks=true,letterpaper=true,citecolor=citecolor,colorlinks,bookmarks=false]{hyperref}
\usepackage[capitalise]{cleveref}
\usepackage[labelfont={bf,small},font=small]{caption}
%\setlength\titlebox{5cm}
\newcommand{\minisection}[1]{\noindent{\textbf{#1.}}}
% You can expand the titlebox if you need extra space
% to show all the authors. Please do not make the titlebox
% smaller than 5cm (the original size); we will check this
% in the camera-ready version and ask you to change it back.


\title{NLP 2024-2025 Homework/Homeworkp LaTeX template}

\author{First Author \\
  Affiliation / Address line 1 \\
  Affiliation / Address line 2 \\
  Affiliation / Address line 3 \\
  {\tt email@domain} \\\And
  Second Author \\
  Affiliation / Address line 1 \\
  Affiliation / Address line 2 \\
  Affiliation / Address line 3 \\
  {\tt email@domain} \\}

\date{}

\begin{document}
\maketitle
\begin{abstract}
 Write here the abstract of your Homework/Homeworkp.
 How to write an abstract? We suggest you to read this~\cite{howabstract}
\end{abstract}

%Describe the NLP task
%Provide examples of input and expected output
%What is the current state-of-the-art for this task (study related work or surveys)
%Applications for this task
%List datasets available for
%Existing tools or libraries that can be used
%Papers with code
%State-of-the-art evaluation


\section{Task description/Problem statement}
Describe in this Section the NLP task addressed in your Homework or Homeworkp.
Provide definitions and cite the proper references.
To correctly cite bibliographical references, you can read this~\cite{latexguide}. This \cite{DevlinCLT19} is an example of bibliographic reference.

\subsection{Examples}
In this Section, provide examples of input and expected output.
Read this~\cite{latexfig}, for some instructions on how to include figures and tables in a LaTeX document


\subsection{Real-world applications}
Describe real-world downstream applications involving the addressed algorithms/models.

\section{Related work}
What is the current state-of-the-art for this task (primary related studies or surveys).
How to write a good Related Work? Read this~\cite{relatedwork}.

\section{Datasets and benchmarks}
Describe available datasets and benchmarks.

\section{Existing tools, libraries, papers with code}
Describe existing tools, libraries that can be used, and papers with code.

\section{State-of-the-art evaluation}
Describe how existing systems/models are evaluated in the literature.

\section{Comparative evaluation}
For Homeworkp only, describe your comparative evaluation experimental setting:
\begin{itemize}
    \item describe the datasets;
    \item describe the systems and the baselines;
    \item describe the measures and the evaluation protocol.
\end{itemize}

\subsection{Results}
Provide evidence about the resulting performances. Discuss the results, and try to motivate the performances. 

\subsection{Discussion}
What are the limits of the compared systems? What are the committed errors? Why?

\section{Conclusions}
Summarize your work, and more important, discuss future works.

\minisection{Contributions} A contribution statement declaring which author did what, in case more than one author.


\bibliographystyle{plain} % We choose the "plain" reference style
\bibliography{refs} % Entries are in the refs.bib file
\end{document}